\chapter{The spherical truncated flow}
\label{chap:sphere_flow}

% Split from the original Gluck_Sphere.tex (stage sections preserved verbatim;
% all labels and \lean pins unchanged).

\paragraph{Formalization map.}
This chapter corresponds primarily to \texttt{Gluck/Sphere/Flow.lean}, \texttt{Gluck/Sphere/StepReparam.lean} and \texttt{Gluck/Sphere/ArcAlgebra.lean}.

\section{Truncated flow layer (S2-B)}
\label{sec:sphere_flow}

\emph{Design decision (iter-046).} The reconstruction field is truncated
\emph{algebraically} — the norm clamped in the numerator, the denominator clamped from
below — so that it becomes globally defined, bounded, and globally Lipschitz in $z$. All
flow machinery (existence on $[0,2\pi]$, uniqueness, continuous dependence, the endpoint
map) is then \emph{unconditional}: no confinement lemma is needed to run the degree
argument. Admissibility is re-imposed a posteriori, only on the single closed trajectory
the winding argument produces (\cref{sec:sphere_admissible}). This removes the
continuation/stitching bookkeeping entirely: one Picard--Lindel\"of application on a
large ball covers $[0,2\pi]$.

\begin{definition}

\leanok
[truncated gauge speed]
  \label{def:truncated_speed}
  \lean{Gluck.truncatedSpeed}
  \uses{def:spherical_speed}
  For parameters $R,\delta$ set
  \[
    \hat q_{\kappa,R,\delta}(\theta,z)\;=\;
      \frac{1+\min(\|z\|,R)^2}
        {2\,\max\bigl(\kappa(\theta)-\langle z,\,i\,e^{i\theta}\rangle_{\mathbb R},\ \delta\bigr)}.
  \]
  On the \emph{admissible set}
  $\{(\theta,z):\|z\|\le R\ \wedge\ \kappa(\theta)-\langle z,i e^{i\theta}\rangle\ge\delta\}$
  both clamps are inactive and $\hat q=q_\kappa$; off it, $\hat q$ is a globally tame
  surrogate.
\end{definition}

\begin{lemma}

\leanok
[truncated speed agrees on the admissible set]
  \label{lem:truncated_speed_eq}
  \lean{Gluck.truncatedSpeed_eq}
  \uses{def:truncated_speed, def:spherical_speed}
  If $\|z\|\le R$ and $\kappa(\theta)-\langle z,i e^{i\theta}\rangle_{\mathbb R}\ge\delta$,
  then $\hat q_{\kappa,R,\delta}(\theta,z)=q_\kappa(\theta,z)$.
\end{lemma}

\begin{proof}

\leanok
  $\min(\|z\|,R)=\|z\|$ by the first hypothesis and
  $\max(\kappa(\theta)-\langle z,i e^{i\theta}\rangle,\delta)
   =\kappa(\theta)-\langle z,i e^{i\theta}\rangle$ by the second; the two formulas
  coincide.
\end{proof}

\begin{lemma}

\leanok
[truncated speed is positive]
  \label{lem:truncated_speed_pos}
  \lean{Gluck.truncatedSpeed_pos}
  \uses{def:truncated_speed}
  If $\delta>0$ then $\hat q_{\kappa,R,\delta}(\theta,z)>0$ for all $\theta,z$.
\end{lemma}

\begin{proof}

\leanok
  The numerator is $\ge 1>0$ (a square plus one) and the denominator is
  $\ge 2\delta>0$; a quotient of positives is positive.
\end{proof}

\begin{lemma}

\leanok
[truncated speed is bounded]
  \label{lem:truncated_speed_le}
  \lean{Gluck.truncatedSpeed_le}
  \uses{def:truncated_speed}
  If $0\le R$ and $0<\delta$ then
  $\hat q_{\kappa,R,\delta}(\theta,z)\le B:=\dfrac{1+R^2}{2\delta}$ for all $\theta,z$.
\end{lemma}

\begin{proof}

\leanok
  $\min(\|z\|,R)\le R$ bounds the numerator by $1+R^2$, and the denominator is
  $\ge 2\delta$; divide.
\end{proof}

\begin{lemma}

\leanok
[truncated speed is Lipschitz in $z$, uniformly in $\theta$]
  \label{lem:truncated_speed_lipschitz}
  \lean{Gluck.truncatedSpeed_lipschitz}
  \uses{def:truncated_speed}
  If $0\le R$ and $0<\delta$ there is $L\ge 0$ (explicitly
  $L=\tfrac{R}{\delta}+\tfrac{1+R^2}{2\delta^2}$ works) such that for every $\theta$ the
  map $z\mapsto\hat q_{\kappa,R,\delta}(\theta,z)$ is $L$-Lipschitz on all of $\mathbb C$.
\end{lemma}

\begin{proof}

\leanok
  Write $\hat q=N/D$ with $N(z)=1+\min(\|z\|,R)^2$ and
  $D(z)=2\max(\kappa(\theta)-\langle z,i e^{i\theta}\rangle,\delta)$. The norm is
  $1$-Lipschitz, the clamp $x\mapsto\min(x,R)$ is $1$-Lipschitz and bounded by $R$ (for
  $x\ge0$), so $N$ is $2R$-Lipschitz (difference of squares of clamped values, both
  $\le R$) and $1\le N\le 1+R^2$. The linear functional
  $z\mapsto\langle z,i e^{i\theta}\rangle$ is $1$-Lipschitz (Cauchy--Schwarz against a unit
  vector) and $x\mapsto\max(x,\delta)$ is $1$-Lipschitz, so $D$ is $2$-Lipschitz and
  $D\ge 2\delta$. Then
  \[
    \Bigl|\frac{N(z)}{D(z)}-\frac{N(w)}{D(w)}\Bigr|
    \le\frac{|N(z)-N(w)|}{D(w)}+N(z)\,\frac{|D(z)-D(w)|}{D(z)D(w)}
    \le\Bigl(\frac{2R}{2\delta}+\frac{(1+R^2)\cdot2}{4\delta^2}\Bigr)|z-w| .
  \]
  (Formalization route: reciprocal is $\tfrac1{(2\delta)^2}$-Lipschitz on
  $[2\delta,\infty)$ by the mean value inequality; compose and multiply
  bounded-Lipschitz factors.)
\end{proof}

\begin{lemma}

\leanok
[truncated speed is jointly continuous]
  \label{lem:truncated_speed_continuous}
  \lean{Gluck.truncatedSpeed_continuous}
  \uses{def:truncated_speed}
  If $\kappa$ is continuous and $0<\delta$, then
  $(\theta,z)\mapsto\hat q_{\kappa,R,\delta}(\theta,z)$ is continuous on
  $\mathbb R\times\mathbb C$.
\end{lemma}

\begin{proof}

\leanok
  Numerator and denominator are continuous ($\min$, $\max$, the norm, $\kappa$, the inner
  product against the continuous unit field $\theta\mapsto i e^{i\theta}$), and the
  denominator never vanishes (it is $\ge 2\delta>0$), so the quotient is continuous
  everywhere — no slab restriction, unlike \cref{lem:spherical_speed_continuous}.
\end{proof}

\begin{definition}

\leanok
[truncated reconstruction field]
  \label{def:truncated_field}
  \lean{Gluck.truncatedField}
  \uses{def:truncated_speed}
  $F_{\kappa,R,\delta}(\theta,z)=\hat q_{\kappa,R,\delta}(\theta,z)\,e^{i\theta}
   \in\mathbb C$. Since $\|e^{i\theta}\|=1$, the field inherits the bound
  $\|F\|\le B$ of \cref{lem:truncated_speed_le} and the Lipschitz constant $L$ of
  \cref{lem:truncated_speed_lipschitz} (the difference at fixed $\theta$ is
  $(\hat q(z)-\hat q(w))\,e^{i\theta}$).
\end{definition}

\begin{lemma}\leanok
[field norm equals speed]
  \label{lem:norm_truncated_field}
  \lean{Gluck.norm_truncatedField}
  \uses{def:truncated_field, lem:truncated_speed_pos}
  For $0<\delta$,
  $\|F_{\kappa,R,\delta}(\theta,z)\|=\hat q_{\kappa,R,\delta}(\theta,z)$ for all
  $\theta,z$.
\end{lemma}

\begin{proof}

\leanok
  $\|\hat q\,e^{i\theta}\|=|\hat q|\,\|e^{i\theta}\|=|\hat q|=\hat q$, since the unit
  field has norm one and $\hat q>0$ by \cref{lem:truncated_speed_pos}.
\end{proof}

\begin{lemma}\leanok
[field Lipschitz bound]
  \label{lem:truncated_field_lipschitz}
  \lean{Gluck.truncatedField_lipschitz}
  \uses{def:truncated_field, lem:truncated_speed_lipschitz}
  For $0\le R$, $0<\delta$ there is a constant $L$ (the one of
  \cref{lem:truncated_speed_lipschitz}) such that for every $\theta$ the map
  $z\mapsto F_{\kappa,R,\delta}(\theta,z)$ is $L$-Lipschitz.
\end{lemma}

\begin{proof}

\leanok
  At fixed $\theta$,
  $F(\theta,z)-F(\theta,w)=(\hat q(\theta,z)-\hat q(\theta,w))\,e^{i\theta}$, whose norm
  is $|\hat q(\theta,z)-\hat q(\theta,w)|\le L\|z-w\|$ by
  \cref{lem:truncated_speed_lipschitz}.
\end{proof}

\begin{lemma}\leanok
[speed Lipschitz bound, uniform in the curvature]
  \label{lem:truncated_speed_lipschitz_uniform}
  % Lean: private `truncatedSpeed_lipschitz_uniform` in Gluck/Sphere/EndpointWinding.lean (private declarations are not pinned)
  \uses{def:truncated_speed, lem:truncated_speed_lipschitz}
  For $0\le R$, $0<\delta$ there is a single $L\ge0$ such that for \emph{every}
  curvature function $\kappa$ and every $\theta$ the map
  $z\mapsto\hat q_{\kappa,R,\delta}(\theta,z)$ is $L$-Lipschitz. (The explicit
  constant of \cref{lem:truncated_speed_lipschitz} depends only on $(R,\delta)$,
  never on $\kappa$; this uniform form breaks the quantifier circularity of the
  winding assembly, where the $L^1$ tolerance $\varepsilon$ must be fixed
  \emph{before} the reparametrized curvature $\kappa\circ h_1$ exists. Landed
  \texttt{private} in stage 1; stage 2 de-privatizes it so the mixed-sign
  assembly in \cref{chap:sphere_mixed} can consume the same witness.)
\end{lemma}

\begin{proof}
  \uses{lem:truncated_speed_lipschitz}
  \leanok
  The proof of \cref{lem:truncated_speed_lipschitz} verbatim, with the constant
  $L=\tfrac{R}{\delta}+\tfrac{1+R^2}{2\delta^2}$ quantified before $\kappa$:
  no step of the numerator/denominator estimate inspects $\kappa$ beyond the
  clamp $\max(\cdot,\delta)$, which is $1$-Lipschitz regardless of its argument.
\end{proof}

\begin{lemma}\leanok
[field Lipschitz bound, uniform in the curvature]
  \label{lem:truncated_field_lipschitz_uniform}
  \lean{Gluck.truncatedField_lipschitz_uniform}
  \uses{def:truncated_field, lem:truncated_speed_lipschitz_uniform}
  For $0\le R$, $0<\delta$ there is a single $L\ge0$ such that for every
  curvature function $\kappa$ and every $\theta$ the map
  $z\mapsto F_{\kappa,R,\delta}(\theta,z)$ is $L$-Lipschitz. (Same
  de-privatization note as \cref{lem:truncated_speed_lipschitz_uniform}.)
\end{lemma}

\begin{proof}
  \uses{lem:truncated_speed_lipschitz_uniform}
  \leanok
  The witness of \cref{lem:truncated_speed_lipschitz_uniform} works: the unit
  frame factor $e^{i\theta}$ has norm one, as in
  \cref{lem:truncated_field_lipschitz}.
\end{proof}

\begin{lemma}\leanok
[field joint continuity]
  \label{lem:truncated_field_continuous}
  \lean{Gluck.truncatedField_continuous}
  \uses{def:truncated_field, lem:truncated_speed_continuous}
  If $\kappa$ is continuous and $0<\delta$, then
  $(\theta,z)\mapsto F_{\kappa,R,\delta}(\theta,z)$ is continuous on
  $\mathbb R\times\mathbb C$.
\end{lemma}

\begin{proof}

\leanok
  Scalar multiplication of the jointly continuous truncated speed
  (\cref{lem:truncated_speed_continuous}) by the continuous unit field
  $\theta\mapsto e^{i\theta}$.
\end{proof}

\begin{lemma}

\leanok
[Picard--Lindel\"of package for the truncated field]
  \label{lem:truncated_field_picard}
  \lean{Gluck.truncatedField_isPicardLindelof}
  \uses{def:truncated_field, lem:truncated_speed_le, lem:truncated_speed_lipschitz,
        lem:truncated_speed_continuous, lem:norm_truncated_field,
        lem:truncated_field_lipschitz, lem:truncated_field_continuous}
  Let $\kappa$ be continuous, $0\le R$, $0<\delta$, and fix an initial-disk radius
  $r_0\ge0$. With $a=r_0+2\pi B+1$, the field $F_{\kappa,R,\delta}$ satisfies the
  Picard--Lindel\"of hypothesis package on the time interval $[0,2\pi]$ with initial time
  $t_0=0$, center $x_0=0$, inner radius $r_0$ and outer radius $a$: it is Lipschitz in the
  space variable with a constant uniform in $\theta$, jointly continuous, norm-bounded by
  $B$, and satisfies the budget condition $B\cdot 2\pi\le a-r_0$.
\end{lemma}

\begin{proof}

\leanok
  All four items are \cref{lem:truncated_speed_lipschitz},
  \cref{lem:truncated_speed_continuous}, \cref{lem:truncated_speed_le}, and the arithmetic
  choice of $a$ respectively. Because the truncated field is bounded and Lipschitz on all
  of $\mathbb C$, restriction to the working ball loses nothing and the budget can always
  be met by enlarging $a$ — this is the payoff of truncation: one application covers
  $[0,2\pi]$ with no continuation argument.
\end{proof}

\begin{lemma}

\leanok
[global flow with continuous dependence]
  \label{lem:spherical_flow_exists}
  \lean{Gluck.exists_sphericalFlow}
  \uses{lem:truncated_field_picard}
  Under the hypotheses of \cref{lem:truncated_field_picard} there is a map
  $\alpha:\mathbb C\times\mathbb R\to\mathbb C$ such that for every $z_0$ in the closed
  disk $\|z_0\|\le r_0$: $\alpha(z_0,0)=z_0$; $\theta\mapsto\alpha(z_0,\theta)$ has on
  $[0,2\pi]$ the derivative $F_{\kappa,R,\delta}(\theta,\alpha(z_0,\theta))$ (within the
  interval); and $\alpha$ is continuous on
  $\{\|z_0\|\le r_0\}\times[0,2\pi]$.
\end{lemma}

\begin{proof}

\leanok
  This is the flow form of the Picard--Lindel\"of theorem
  (in Mathlib:
  \texttt{IsPicardLindelof.exists\_forall\_mem\_closedBall\_eq\_hasDerivWithinAt\_continuousOn}),
  applied to the package of \cref{lem:truncated_field_picard}. It returns a single
  $\alpha$ defined on the product with exactly the three listed properties.
\end{proof}

\begin{definition}

\leanok
[spherical flow]
  \label{def:spherical_flow}
  \lean{Gluck.sphericalFlow}
  \uses{lem:spherical_flow_exists}
  Fix, once per parameter tuple $(\kappa,R,\delta,r_0)$, a choice
  $\Phi=\Phi_{\kappa,R,\delta,r_0}:\mathbb C\times\mathbb R\to\mathbb C$ of the map
  supplied by \cref{lem:spherical_flow_exists}. (The whole flow is chosen at once — not a
  separate choice per initial point — so downstream continuity statements can consume it.)
\end{definition}

\begin{lemma}

\leanok
[flow specification]
  \label{lem:spherical_flow_spec}
  \lean{Gluck.sphericalFlow_spec}
  \uses{def:spherical_flow}
  For $\|z_0\|\le r_0$: $\Phi(z_0,0)=z_0$ and $\theta\mapsto\Phi(z_0,\theta)$ solves
  $\gamma'=F_{\kappa,R,\delta}(\theta,\gamma)$ on $[0,2\pi]$ (derivative within the interval).
\end{lemma}

\begin{proof}

\leanok
  Unfold the choice of \cref{def:spherical_flow}.
\end{proof}

\begin{lemma}

\leanok
[flow continuity]
  \label{lem:spherical_flow_continuousOn}
  \lean{Gluck.sphericalFlow_continuousOn}
  \uses{def:spherical_flow}
  $\Phi$ is continuous on $\{\|z_0\|\le r_0\}\times[0,2\pi]$.
\end{lemma}

\begin{proof}

\leanok
  Unfold the choice of \cref{def:spherical_flow}.
\end{proof}

\begin{lemma}

\leanok
[flow uniqueness]
  \label{lem:spherical_flow_unique}
  \lean{Gluck.sphericalFlow_unique}
  \uses{def:spherical_flow, lem:truncated_speed_lipschitz, lem:spherical_flow_spec}
  If $g:[0,2\pi]\to\mathbb C$ solves $g'=F_{\kappa,R,\delta}(\theta,g(\theta))$ on
  $[0,2\pi]$ with $g(0)=z_0$, $\|z_0\|\le r_0$, then $g=\Phi(z_0,\cdot)$ on $[0,2\pi]$.
\end{lemma}

\begin{proof}

\leanok
  Both solve the same ODE with the same initial value; the field is Lipschitz in the space
  variable uniformly in time (\cref{lem:truncated_speed_lipschitz}), so the standard
  uniqueness theorem for ODEs (Gr\"onwall; in Mathlib
  \texttt{ODE\_solution\_unique\_of\_mem\_Icc}) applies. Uniqueness is what later
  identifies explicitly constructed trajectories — circular arcs, reflected trajectories —
  with the flow.
\end{proof}

\begin{lemma}\leanok
[two-solution uniqueness on a subinterval]
  \label{lem:truncated_field_solution_unique}
  \lean{Gluck.truncatedField_solution_unique}
  \uses{def:truncated_field, lem:truncated_field_lipschitz}
  Let $0\le R$, $0<\delta$ and $0\le T$. If $g_1,g_2:[0,T]\to\mathbb C$ both solve
  $g'=F_{\kappa,R,\delta}(\theta,g(\theta))$ on $[0,T]$ and $g_1(0)=g_2(0)$, then
  $g_1=g_2$ on $[0,T]$. (Unlike \cref{lem:spherical_flow_unique} this compares two
  arbitrary solutions, on an arbitrary compact interval $[0,T]$ — not necessarily
  $[0,2\pi]$ — and makes no reference to the chosen flow $\Phi$; the field is defined
  and Lipschitz for all times, so no time restriction is needed.)
\end{lemma}

\begin{proof}

\leanok
  The field is Lipschitz in the space variable uniformly in time
  (\cref{lem:truncated_field_lipschitz}) and each $g_i$ is continuous on $[0,T]$ (it has
  a derivative there), so the standard two-solution uniqueness theorem applies on
  $[0,T]$ with initial time at the left endpoint — the same Gr\"onwall-type argument as
  \cref{lem:spherical_flow_unique}, with $2\pi$ replaced by $T$.
\end{proof}

\begin{definition}

\leanok
[spherical endpoint map]
  \label{def:spherical_endpoint}
  \lean{Gluck.sphericalEndpoint}
  \uses{def:spherical_flow}
  $E(z_0)=\Phi(z_0,2\pi)-z_0$. A zero of $E$ is a closed trajectory of the truncated
  flow.
\end{definition}

\begin{lemma}

\leanok
[endpoint map continuity]
  \label{lem:spherical_endpoint_continuousOn}
  \lean{Gluck.sphericalEndpoint_continuousOn}
  \uses{def:spherical_endpoint, lem:spherical_flow_continuousOn}
  $E$ is continuous on the closed disk $\{\|z_0\|\le r_0\}$.
\end{lemma}

\begin{proof}

\leanok
  $z_0\mapsto\Phi(z_0,2\pi)$ is continuous by \cref{lem:spherical_flow_continuousOn}
  (restriction of a jointly continuous map to a time slice), and subtracting the identity
  preserves continuity.
\end{proof}

